\appendix
\section{Artifact Appendix}

\paragraph{Artifact and badge scope.}
SCARF applies for Artifact Available, Artifacts Evaluated--Functional, and
Results Reproduced. The release is identified by Zenodo DOI
\url{https://doi.org/10.5281/zenodo.21482385}. It contains the functional and
cycle simulator, FSDR and SAES, three model adapters, Chisel RTL, experiment
automation, result validators, a CUDA container, and a public physical-design
proxy. The machine-readable catalog maps every Evaluation figure and table to
its command, raw inputs, evidence class, and acceptance gate.

\paragraph{Container workflow.}
The Docker image uses CUDA~12.1 and the pinned Python~3.10 classic profile.
Its default command downloads the hash-pinned MVSplat quick checkpoint, runs
the redistributable synthetic fixture, and writes a structured Functional
record. The full quality and mechanism workflows use the same record schema
with the official prepared datasets and checkpoints.

\paragraph{Key results and hardware.}
Figure~8 uses a Jetson Orin NX 16\,GB in MAXN mode with the GPU locked at
918\,MHz. The measurement wrapper records CUDA event repetitions, Nsight
Systems output, \texttt{tegrastats}, JetPack/L4T/CUDA/PyTorch versions, clock
state, power mode, and temperature. The archive also provides the normalized
nine-pair Figure~8 reference table for reviewers without an Orin NX. Reference
results never substitute for generated raw evidence. Full model evaluation
uses an NVIDIA GPU with at least 24\,GB VRAM; quick Functional inference uses
an 8\,GB CUDA GPU. CPU checks cover schemas and scaling. RTL uses JDK~11+,
sbt~1.9+, Chisel~6.6, and Verilator~5+.

\paragraph{Data and protocol.}
Re10K, ACID, and official checkpoints are downloaded from manifest-pinned
sources. DL3DV-Benchmark is auto-gated and each evaluator accepts the upstream
terms independently; the dataset is not redistributed in Zenodo. The protocol
preserves the committed upstream evaluation indices: 6,474 executable Re10K
samples, 1,595 ACID samples, and 140 DL3DV samples with exact context and
target views. Each result binds source index, ordered selection, prepared tree,
checkpoint, runtime assets, environment, command, and output by SHA256.

\paragraph{Calibration and execution.}
The global mechanism configuration follows a fixed contract with 24 DL3DV
calibration scenes and eight disjoint DL3DV holdout scenes. Configuration,
route decisions, retained descriptors, Gaussian materialization, FSDR events,
and cycle traces are bound into each result. Quality records retain PSNR, SSIM,
LPIPS, and per-target-view aggregation; mechanism records retain cache,
candidate, buffer, tile, Gaussian, and MMCU-slot events.

\paragraph{Public hardware proxy.}
The public hardware path routes the released RTL with iFlow and ASAP7, then
applies the paper-cited DeepScaleTool normalization to generate a labeled
28\,nm-equivalent estimate. Generated reports preserve ASAP7 raw values,
scaling factors, estimates, and manuscript comparison targets as separate
fields. The commercial TSMC 28\,nm PDK, memory compiler, Liberty files, and
LPDDR PHY are not redistributed.

\paragraph{Acceptance and variation.}
Quality tolerances are 0.15\,dB PSNR, 0.005 SSIM, and 0.005 LPIPS. Figure~8
and Figure~11 geometric means allow 5\% relative variation under their fixed
contracts. Mechanism rates and utilization allow two percentage points; event
counts allow 5\% relative variation. Deterministic schemas, hashes, selection
identities, RTL, and scaling formulas must match exactly.
